\section*{Bab \ref{ch:b2:measurement-statistics} --- Statistika Pengukuran Kimia}

\begin{solution}{exo:b2:measurement-statistics:1}
$\bar V = 62.41/5 = \qty{12.482}{mL}$; $\sum(V_i - \bar V)^2 = 0.00328$, $s = \sqrt{0.00328/4} \approx
\qty{0.029}{mL}$; $s/\sqrt5 \approx \qty{0.013}{mL}$.
\end{solution}

\begin{solution}{exo:b2:measurement-statistics:2}
$t = 2.776$ untuk 4 derajat kebebasan: $12.482 \pm 2.776 \times 0.0128$, yaitu $(12.48 \pm 0.04)$~mL.
\end{solution}

\begin{solution}{exo:b2:measurement-statistics:3}
Ketidakpastian relatif massa dan volume adalah 0.04~\% dan 0.06~\%; jika
digabungkan secara
kuadratur,
\[\frac{u(c)}{c} = \sqrt{\Bigl(\frac{0.0002}{0.5000}\Bigr)^2 + \Bigl(\frac{0.15}{250.0}\Bigr)^2}
\approx 0.07~\%.\]
\end{solution}

\begin{solution}{exo:b2:measurement-statistics:4}
$\bar x = 1.5$, $\bar y = 0.31$, $S_{xx} = 5$, $S_{xy} = 0.99$: $b = 0.198$, $a = 0.31 - 0.198 \times 1.5
= 0.013$.
\end{solution}

\begin{solution}{exo:b2:measurement-statistics:5}
$|10.12 - 10.00|/(0.08/\sqrt5) = 0.12/0.036 \approx 3.4 > 2.776$: perbedaannya
bermakna, metode itu \omterm{def:b2:measurement-statistics:validation}{bias} (sekitar \qty{+0.12}{mg/L}).
\end{solution}

\begin{solution}{exo:b2:measurement-statistics:6}
$\mathrm{pH} = -\log_{10}c$, $\mathrm d\,\mathrm{pH}/\mathrm dc = -1/(c\ln10)$:
$u(\mathrm{pH}) = (u(c)/c)/\ln10 = 0.02/2.303 \approx 0.009$.
\end{solution}

\begin{solution}{exo:b2:measurement-statistics:7}
$x_{\mathrm{LOD}} = 3 \times 0.0012/0.0098 \approx \qty{0.37}{\micro g/L}$; $x_{\mathrm{LOQ}} = 10
\times 0.0012/0.0098 \approx \qty{1.2}{\micro g/L}$.
\end{solution}

\begin{solution}{exo:b2:measurement-statistics:8}
Garisnya memberikan \qty{3.11}{mg/L} dalam larutan yang diukur; sampelnya
telah diencerkan lima kali:
$3.11 \times 50.0/10.0 \approx \qty{15.6}{mg/L}$.
\end{solution}

\begin{solution}{exo:b2:measurement-statistics:9}
\qty{0.05}{mg/L} adalah \omterm{def:b2:measurement-statistics:validation}{keterulangan}, \qty{0.15}{mg/L} adalah \omterm{def:b2:measurement-statistics:validation}{ketertiruan}.
Setiap laboratorium mempunyai \omterm{def:b2:measurement-statistics:validation}{bias} kecilnya sendiri (kalibrasi, standar,
instrumen); \omterm{def:b2:measurement-statistics:validation}{bias-bias} ini berbeda antarlaboratorium dan menambah sebaran,
sehingga \omterm{def:b2:measurement-statistics:validation}{ketertiruan} selalu lebih besar.
\end{solution}

\begin{solution}{exo:b2:measurement-statistics:10}
$\partial S/\partial a = 0$ berbunyi $\sum(y_i - a - bx_i) = 0$: jumlah \omterm{def:b2:measurement-statistics:calibration}{sisaan-sisaannya} nol. Dengan membagi dengan
$n$: $\bar y = a + b\bar x$, sehingga $(\bar x, \bar y)$ terletak pada garis itu.
\end{solution}

\begin{solution}{exo:b2:measurement-statistics:11}
Satu tahap:
\[\sqrt{(0.006/1.000)^2 + (0.08/100.0)^2} \approx 0.61~\%.\]
Satu tahap sepuluh kali:
\[\sqrt{(0.02/10.00)^2 + (0.08/100.0)^2} \approx 0.22~\%,\]
dan dua tahap seperti itu yang digabungkan secara kuadratur memberikan
$0.22\sqrt2 \approx 0.30~\%$. Jalur dua tahap dua kali lebih presisi: pipet
kecil mendominasi jalur satu tahap.
\end{solution}

\begin{solution}{exo:b2:measurement-statistics:12}
$z = 1.6/\sqrt{0.4^2 + 0.5^2} \approx 2.5 > 2$: tidak saling sesuai;
sekurang-kurangnya satu laboratorium mempunyai \omterm{def:b2:measurement-statistics:validation}{bias}. Dengan ketidakpastian
diperluas ($k = 2$), yang pertama memberikan $9.6 \pm 0.8$, yang memuat 10;
yang kedua $11.2 \pm 1.0$, seluruhnya di atas 10. Sebelum keputusan apa pun
tentang air itu, ketidaksesuaian itu harus diselesaikan, misalnya dengan
bahan acuan.
% ledger: who:lead
\end{solution}

\begin{solution}{pb:b2:measurement-statistics:1}
\textbf{1.} Asam mengubah atomisasi dan tanggapannya; standar dan sampel
harus mempunyai matriks yang sama agar kemiringannya berlaku.
\textbf{2.} $\bar x = 10$, $\bar y = 0.1004$, $S_{xx} = 250$, $S_{xy} = 2.445$: $b = 0.00978$ per
\unit{\micro g/L}, $a = 0.1004 - 0.0978 = 0.0026$.
\textbf{3.} $+0.0004$, $-0.0005$, $+0.0006$, $-0.0013$, $+0.0008$.
\textbf{4.} Tandanya berselang-seling, tanpa kelengkungan: garis itu memadai
pada rentang 0--\qty{20}{\micro g/L}.
\textbf{5.} $s_{y/x} = \sqrt{3.1 \times 10^{-6}/3} \approx 0.00102$.
\textbf{6.} Blangko (asam, air, tanur) memberikan absorbans kecilnya sendiri.
\textbf{7.} Kepekaan: absorbans per \unit{\micro g/L} timbal.
\textbf{8.} $\bar y_0 = 0.1010$; $s = 0.0030$.
\textbf{9.} Konsentrasinya $x_0 = (0.1010 - 0.0026)/0.00978 \approx \qty{10.06}{\micro g/L}$.
\textbf{10.} Simpangan bakunya adalah $s_{x_0} = (0.00102/0.00978)\sqrt{1/3 + 1/5 + (0.1010 - 0.1004)^2/(0.00978^2 \times 250)}
\approx 0.104 \times 0.730 \approx \qty{0.076}{\micro g/L}$.
\textbf{11.} $n - 2 = 3$; $t = 3.182$.
\textbf{12.} $10.06 \pm 3.182 \times 0.076 = (10.06 \pm 0.24)$~\unit{\micro g/L}.
\textbf{13.} Suku $1/m$ turun dari 1 menjadi $1/3$: $s_{x_0}$ berkurang sekitar sepertiga.
\textbf{14.} $f = 50.50/50.0 = 1.010$; $10.06 \times 1.010 \approx \qty{10.16}{\micro g/L}$.
\textbf{15.} $f = 1 + V_2/V_1$: $u^2(f) = (u(V_2)/V_1)^2 + (V_2u(V_1)/V_1^2)^2 = (0.005/50.0)^2 +
(0.50 \times 0.03/2500)^2$, $u(f) \approx 1.0 \times 10^{-4}$, secara relatif 0.01~\%.
\textbf{16.} Tidak: 0.01~\% dibandingkan dengan $0.076/10.06 \approx 0.75~\%$.
\textbf{17.} $s_{\mathrm{blank}} \approx 0.00115$; $x_{\mathrm{LOD}} = 3 \times 0.00115/0.00978 \approx
\qty{0.35}{\micro g/L}$.
\textbf{18.} $10 \times 0.00115/0.00978 \approx \qty{1.2}{\micro g/L}$.
\textbf{19.} Ya, jauh di atasnya.
\textbf{20.} Ya: dari 9.92 sampai \qty{10.40}{\micro g/L} dalam air keran.
\textbf{21.} Satu pembacaan saja memberikan $(0.104 - 0.0026)/0.00978 \times 1.010 \approx \qty{10.5}{\micro g/L}$,
yang jika dibaca sendiri mengisyaratkan kesimpulan yakin ``di atas
pedoman'' yang tidak didukung oleh data.
\textbf{22.} Aturan yang dipakai untuk membangun selang itu menangkap nilai
sebenarnya pada 95~\% deret yang padanya aturan itu diterapkan; selang itu
bukan peluang tentang satu nilai ini.
\textbf{23.} Lebih banyak pembacaan ulangan dan lebih banyak standar di
dekat \qty{10}{\micro g/L} (suku $1/m$ dan
$1/n$), atau metode yang lebih peka; selang itu menyempit hanya perlahan,
sebanding dengan $1/\sqrt{\text{banyaknya}}$.
\textbf{24.} Dengan menganalisis air acuan bersertifikat, atau dengan
menambahkan sejumlah timbal yang diketahui ke dalam sampel dan memeriksa
perolehan kembalinya.
\textbf{25.} Di dekat \qty{10}{\micro g/L}, ketidakpastian metode rutin
adalah beberapa persen, seperti di sini: nilai pedoman yang lebih rendah
tidak dapat diperiksa secara andal oleh laboratorium biasa.
\textbf{26.} $\boldsymbol{(10.16 \pm 0.24)}$~\unit{\micro g/L} pada 95~\%:
selang itu memuat \qty{10}{\micro g/L}, sehingga pengukuran itu tidak
menunjukkan bahwa air itu melampaui pedoman, dan tidak pula bahwa air itu
memenuhinya dengan selisih yang aman.
\end{solution}
